\documentclass[12pt,a4paper]{article}
\usepackage{fontspec}
\IfFontExistsTF{Times New Roman}{\setmainfont{Times New Roman}}{\IfFontExistsTF{TeX Gyre Termes}{\setmainfont{TeX Gyre Termes}}{\setmainfont{Tinos}}}
\IfFontExistsTF{Consolas}{\setmonofont{Consolas}[Scale=0.8]}{\setmonofont{Liberation Mono}[Scale=0.73]}
\usepackage[russian]{babel}
\usepackage[left=30mm,right=15mm,top=20mm,bottom=20mm]{geometry}
\usepackage{setspace}\onehalfspacing
\usepackage{tikz}\usetikzlibrary{arrows.meta,positioning,calc}
\usepackage{listings,xcolor,float,hyperref}
\usepackage{enumitem}
\usepackage{titlesec}
\titlespacing*{\section}{0pt}{2.1ex plus 0.2ex}{1ex}
\titlespacing*{\subsection}{0pt}{1.4ex plus 0.2ex}{0.7ex}
\hypersetup{hidelinks,breaklinks=true,
  pdftitle={Отчёт по лабораторной работе № 5. Динамические библиотеки},
  pdfauthor={Прудников Я. А.}}
\setlength{\parindent}{1.25cm}
\emergencystretch=2em
\definecolor{codegray}{RGB}{247,247,247}
\lstset{language=C,basicstyle=\ttfamily\footnotesize\singlespacing,
  backgroundcolor=\color{codegray},breaklines=true,showstringspaces=false,
  frame=single,framesep=5pt,columns=fullflexible,keepspaces=true}
\setlist{itemsep=2pt,topsep=4pt}
\begin{document}

\begin{titlepage}
\begin{center}
{\small МИНИСТЕРСТВО НАУКИ И ВЫСШЕГО ОБРАЗОВАНИЯ\\РОССИЙСКОЙ ФЕДЕРАЦИИ}\\[4mm]
{\small Федеральное государственное автономное образовательное учреждение\\ высшего образования\\ «Московский авиационный институт (национальный исследовательский университет)»}
\vfill
{\bfseries\large ОТЧЁТ ПО ЛАБОРАТОРНОЙ РАБОТЕ № 5}\\[9mm]
Дисциплина:\\[2mm]
{\bfseries Операционные системы и архитектура компьютера}\\[6mm]
Тема:\\[2mm]
{\bfseries Динамические библиотеки}\\[3mm]
Контракты № 3 и № 7
\vfill
\end{center}
\begin{flushright}\begin{minipage}{0.45\textwidth}
Выполнил:\\ студент группы М80-203БВ-25\\ Прудников Я. А.\par\vspace{7mm}
Преподаватель:\\ Миронов Е. С.
\end{minipage}\end{flushright}
\vfill\begin{center}Москва, 2026 г.\end{center}
\end{titlepage}

\section{Цель работы}
Приобрести практические навыки создания динамических библиотек и использования их функций двумя способами: посредством линковки с библиотекой при сборке программы и посредством явной загрузки библиотеки во время исполнения средствами операционной системы. Освоить организацию кроссплатформенного интерфейса загрузки библиотек для Linux и Windows.

\section{Постановка задачи}
Разработать две динамические библиотеки с одинаковыми контрактами функций и разными алгоритмами их реализации. Подготовить две тестовые программы: первая должна использовать одну из библиотек при линковке, вторая~- загружать библиотеки по относительным путям и переключаться между ними во время работы.

Для выбранных контрактов заданы следующие функции:
\begin{itemize}
  \item \texttt{int PrimeCount(int A, int B)}~- количество простых чисел на отрезке $[A,B]$, где $A$ и $B$~- натуральные числа. Первая реализация использует перебор делителей, вторая~- решето Эратосфена.
  \item \texttt{float Square(float A, float B)}~- площадь фигуры по двум сторонам. Первая реализация вычисляет площадь прямоугольника $S=AB$, вторая~- площадь прямоугольного треугольника $S=AB/2$.
\end{itemize}

Пользователь вводит команды \texttt{1 A B} и \texttt{2 A B} для вызова соответствующих функций. В программе № 2 команда \texttt{0} переключает активную библиотеку. В программе № 1 переключение не предусмотрено. Сборка и использование библиотек должны быть предусмотрены для Linux и Windows.

\section{Архитектура решения}
В проекте выделены два файла с реализациями контрактов~- \texttt{lib1.c} и \texttt{lib2.c}. При сборке они преобразуются в динамические библиотеки: \texttt{lib1.so} и \texttt{lib2.so} в Linux либо \texttt{lib1.dll} и \texttt{lib2.dll} в Windows. Объявления математических функций находятся в \texttt{contracts.h}.

Программа \texttt{program1} связывается с первой библиотекой при линковке и вызывает \texttt{PrimeCount()} и \texttt{Square()} непосредственно. Программа \texttt{program2} обращается к модулю \texttt{lib\_manager.c}, который хранит дескриптор активной библиотеки, её номер и указатели на функции.

Для кроссплатформенности операции загрузки, получения адреса и освобождения библиотеки вынесены в общий интерфейс \texttt{os\_api.h}. В Linux его реализует файл \texttt{platform/os\_linux.c}, а в Windows~- \texttt{platform/os\_windows.c}. При сборке выбирается только один из этих файлов. Алгоритмы вычислений и обработка команд остаются общими для обеих ОС.

\subsection{Схема взаимодействия компонентов}
\begin{figure}[H]
\centering
\begin{tikzpicture}[
  >=Stealth,
  every node/.style={font=\small},
  box/.style={draw,rounded corners,align=center,minimum width=34mm,minimum height=10mm,fill=black!3},
  node distance=8mm and 8mm
]
\node[box] (p1) {\texttt{program1}\\прямой вызов};
\node[box,right=39mm of p1] (p2) {\texttt{program2}\\обработка команд};
\node[box,below=8mm of p2] (manager) {\texttt{lib\_manager}\\дескриптор и указатели};
\node[box,below=8mm of manager] (api) {\texttt{os\_api.h}\\единый интерфейс};
\node[box,below left=9mm and -10mm of api] (linux) {\texttt{os\_linux.c}\\POSIX};
\node[box,below right=9mm and -10mm of api] (win) {\texttt{os\_windows.c}\\WinAPI};
\node[box,below=17mm of linux,xshift=-17mm] (lib1) {\texttt{lib1.so / lib1.dll}\\перебор / прямоугольник};
\node[box,right=11mm of lib1] (lib2) {\texttt{lib2.so / lib2.dll}\\решето / треугольник};
\draw[->] (p2) -- (manager);
\draw[->] (manager) -- (api);
\draw[->] (api) -- (linux);
\draw[->] (api) -- (win);
\coordinate (join) at ($(linux.south)!0.5!(win.south)+(0,-5mm)$);
\draw[dashed] (linux.south) |- (join);
\draw[dashed] (win.south) |- (join);
\draw[->,dashed] (join) -| (lib1.north);
\draw[->,dashed] (join) -| (lib2.north);
\draw[->] (p1.south) -- ++(0,-8mm) -| ([xshift=-15mm]lib1.north) node[pos=.20, left]{линковка};
\end{tikzpicture}
\caption{Общая архитектура для Linux и Windows}
\end{figure}

\subsection{Файловая структура проекта}
\begin{itemize}
  \item \texttt{src/lib1.c}, \texttt{src/lib2.c}~- две реализации контрактов;
  \item \texttt{src/include/contracts.h}~- общие объявления функций;
  \item \texttt{src/program1.c}~- программа с линковкой библиотеки при сборке;
  \item \texttt{src/program2.c}~- программа с явной загрузкой библиотеки;
  \item \texttt{src/lib\_manager.c}, \texttt{src/include/lib\_manager.h}~- управление текущей библиотекой;
  \item \texttt{src/include/os\_api.h}~- общий интерфейс загрузки;
  \item \texttt{src/platform/os\_linux.c} и \texttt{src/platform/os\_windows.c}~- платформозависимые реализации;
  \item \texttt{CMakeLists.txt}~- описание сборки библиотек и программ.
\end{itemize}

\section{Реализация}
\subsection{Функции библиотек}
В первой реализации \texttt{PrimeCount()} перебирает числа на заданном отрезке и проверяет, существуют ли у каждого числа делители, отличные от единицы и самого числа. Если такие делители отсутствуют, число учитывается в результате. Вторая реализация строит решето Эратосфена: числа до верхней границы помечаются как потенциально простые, после чего вычёркиваются кратные найденных простых чисел. Подсчёт результата выполняется только на отрезке $[A,B]$.

Для функции \texttt{Square()} используются две формулы. В первой библиотеке вычисляется $S=AB$, а во второй~- $S=AB/2$. Совпадение сигнатур позволяет менять реализацию без изменения вызывающего кода.

\subsection{Использование библиотеки при линковке}
Файл \texttt{program1.c} подключает \texttt{contracts.h}. При вводе команды \texttt{1} вызывается \texttt{PrimeCount()}, при вводе \texttt{2}~- \texttt{Square()}. Цель \texttt{lib1} указывается в \texttt{target\_link\_libraries()} при сборке исполняемого файла. Программа не осуществляет поиск символов самостоятельно.

\subsection{Динамическая загрузка и переключение}
Программа \texttt{program2} использует собственный интерфейс из трёх операций: \texttt{lib\_open()}, \texttt{get\_address()} и \texttt{lib\_close()}. В Linux они реализуются через \texttt{dlopen()}, \texttt{dlsym()} и \texttt{dlclose()} соответственно. Для получения адресов используются имена экспортируемых символов \texttt{PrimeCount} и \texttt{Square}.

Изначально выбирается первая библиотека. После успешной загрузки её дескриптор и адреса функций сохраняются в статических переменных модуля \texttt{lib\_manager.c}. При вводе команды \texttt{0} текущая библиотека освобождается, её номер меняется и выполняется загрузка другой библиотеки с обновлением адресов функций. Поэтому одинаковый вызов \texttt{call\_square()} выполняет разные реализации до и после переключения.

При завершении работы программа освобождает дескриптор библиотеки. Предусмотрена проверка результатов загрузки и получения адресов функций; при ошибке инициализации или переключения работа прекращается. После закрытия библиотеки указатели на функции обнуляются.

\subsection{Кроссплатформенность}
В Linux используются функции POSIX из \texttt{dlfcn.h}, в Windows~- функции WinAPI из \texttt{windows.h}. Обе реализации предоставляют одинаковый интерфейс для менеджера библиотек. Соответствие операций представлено в таблице.

\begin{table}[H]
\centering\small
\begin{tabular}{|p{48mm}|p{42mm}|p{50mm}|}
\hline
Операция & Linux & Windows \\
\hline
Загрузка библиотеки & \texttt{dlopen()} & \texttt{LoadLibraryA()} \\
Получение адреса & \texttt{dlsym()} & \texttt{GetProcAddress()} \\
Освобождение & \texttt{dlclose()} & \texttt{FreeLibrary()} \\
Диагностика ошибок & \texttt{dlerror()} & \texttt{GetLastError()} \\
Формат файла & \texttt{.so} & \texttt{.dll} \\
\hline
\end{tabular}
\caption{Сопоставление интерфейсов загрузки библиотек}
\end{table}

В Linux \texttt{dlopen()} с флагом \texttt{RTLD\_NOW} возвращает дескриптор типа \texttt{void *}. В Windows \texttt{LoadLibraryA()} возвращает дескриптор \texttt{HMODULE}. При ошибке загрузки обе функции возвращают \texttt{NULL}. Функция \texttt{GetProcAddress()} возвращает указатель типа \texttt{FARPROC}, тогда как \texttt{dlsym()} возвращает \texttt{void *}. Преобразования адресов в типы указателей на функции выполняются с учётом соглашений целевых платформ.

Функция \texttt{FreeLibrary()} возвращает ненулевое значение при успехе, а \texttt{dlclose()}~- ноль. В общем интерфейсе результат закрытия нормализован: ноль при успехе и $-1$ при ошибке. Относительные пути выбираются макросом \texttt{\_WIN32}.

CMake выбирает платформозависимый файл по условию \texttt{WIN32}. Библиотеки создаются как \texttt{SHARED}; для Windows с MS-совместимыми инструментами сборки включён автоматический экспорт символов. Это позволяет обращаться к функциям из DLL как при линковке, так и по имени во время исполнения. Для Linux сохраняется подключение \texttt{CMAKE\_DL\_LIBS}, необходимое там, где загрузчик требует отдельной библиотеки. Для сборки в Windows и WSL используются разные каталоги, поскольку кэш CMake зависит от среды сборки.

\section{Ключевые фрагменты кода}
Приведены основные фрагменты интерфейсов и логики работы с динамическими библиотеками.

\subsection{Контракты функций}
\begin{lstlisting}
int PrimeCount(int A, int B);
float Square(float A, float B);
\end{lstlisting}
\begin{center}\small Листинг 1 - Контракты из \texttt{contracts.h}\end{center}

\subsection{Хранение состояния менеджера}
\begin{lstlisting}
static void *descriptor = NULL;
static int current_lib = 0;

static int (*prime_count)(int, int) = NULL;
static float (*square)(float, float) = NULL;

#ifdef _WIN32
static const char *paths[] = {
    "./lib1.dll", "./lib2.dll"
};
#else
static const char *paths[] = {
    "./lib1.so", "./lib2.so"
};
#endif
\end{lstlisting}
\begin{center}\small Листинг 2 - Состояние модуля \texttt{lib\_manager.c}\end{center}

\subsection{Получение адресов функций}
\begin{lstlisting}
descriptor = lib_open(paths[current_lib]);
if (descriptor == NULL) {
    return -1;
}

prime_count = (int (*)(int, int))
    get_address(descriptor, "PrimeCount");
square = (float (*)(float, float))
    get_address(descriptor, "Square");
\end{lstlisting}
\begin{center}\small Листинг 3 - Фрагмент инициализации библиотеки\end{center}

\subsection{Вызов функций в двух программах}
\begin{lstlisting}
/* program1.c */
int result_p = PrimeCount(arg1_p, arg2_p);
float result_s = Square(arg1_s, arg2_s);

/* program2.c */
int result_p = call_prime_count(arg1_p, arg2_p);
float result_s = call_square(arg1_s, arg2_s);
\end{lstlisting}
\begin{center}\small Листинг 4 - Различие в способах вызова функций\end{center}

\subsection{Интерфейс загрузки и реализация для Windows}
Общий интерфейс скрывает платформенные типы и способы вызова системных функций.
\begin{lstlisting}
void *lib_open(const char *filename);
void *get_address(void *descriptor, const char *query);
int lib_close(void *descriptor);
\end{lstlisting}
\begin{center}\small Листинг 5 - Интерфейс \texttt{os\_api.h}\end{center}

\begin{lstlisting}
#include "os_api.h"
#include <windows.h>
void *lib_open(const char *filename){
    return LoadLibraryA(filename);
}
void *get_address(void *descriptor, const char *query){
    return (void *)GetProcAddress((HMODULE)descriptor, query);
}
int lib_close(void *descriptor){
    if (FreeLibrary((HMODULE)descriptor) == 0){
        return -1;
    }
    return 0;
}
\end{lstlisting}
\begin{center}\small Листинг 6 - Реализация интерфейса для Windows\end{center}

\subsection{Выбор платформы при сборке}
\begin{lstlisting}[language={}]
add_library(lib1 SHARED src/lib1.c)
add_library(lib2 SHARED src/lib2.c)

set_target_properties(lib1 lib2 PROPERTIES
    PREFIX ""
    WINDOWS_EXPORT_ALL_SYMBOLS ON
)

add_executable(program1 src/program1.c)
target_link_libraries(program1 PRIVATE lib1)

if(WIN32)
    set(OS_SOURCE src/platform/os_windows.c)
else()
    set(OS_SOURCE src/platform/os_linux.c)
endif()

add_executable(program2
    src/program2.c
    src/lib_manager.c
    ${OS_SOURCE}
)
target_link_libraries(program2 PRIVATE ${CMAKE_DL_LIBS})
add_dependencies(program2 lib1 lib2)
\end{lstlisting}
\begin{center}\small Листинг 7 - Основная часть \texttt{CMakeLists.txt}\end{center}

\section{Тестирование}
Функциональная проверка выполнена в Linux (WSL) и Windows. На обеих платформах проведённые тесты завершились успешно. Ниже приведены результаты консольных запусков обеих программ в Linux.

\subsection{Программа № 2: динамическая загрузка}
\begin{lstlisting}[language={},numbers=none]
1 2 4
2
2 2 4
8.000000
0
Library switched
2 2 4
4.000000
0
Library switched
\end{lstlisting}
\begin{center}\small Листинг 8 - Консольный ввод и вывод \texttt{program2}\end{center}

Команда \texttt{1 2 4} вернула значение 2: на отрезке $[2,4]$ простыми являются числа 2 и 3. Для команды \texttt{2 2 4} первая библиотека вычислила площадь прямоугольника $2\cdot4=8$. После ввода \texttt{0} была выбрана вторая библиотека и та же команда вернула площадь прямоугольного треугольника $2\cdot4/2=4$. Повторная команда \texttt{0} переключила реализацию обратно.

\subsection{Программа № 1: подключение при линковке}
\begin{lstlisting}[language={},numbers=none]
1 1 100
25
2 4 6
24.000000
0
Unknown command
\end{lstlisting}
\begin{center}\small Листинг 9 - Консольный ввод и вывод \texttt{program1}\end{center}

На отрезке $[1,100]$ получено 25 простых чисел. Площадь прямоугольника со сторонами 4 и 6 равна 24. Ввод команды \texttt{0} вызвал сообщение \texttt{Unknown command}, поскольку программа № 1 не предусматривает переключение библиотек.

Приведённые примеры подтверждают корректность показанных вычислений и переключения реализации функции \texttt{Square()}. Измерение производительности алгоритмов и расширенное тестирование обработки ошибочного ввода в эти результаты не входят.

\subsection{Сборка и проверка в Windows}
Для сборки под Windows с генератором Visual Studio применяется следующая последовательность команд PowerShell:
\begin{lstlisting}[language={}]
cmake -S . -B build-win
cmake --build build-win --config Debug
cd build-win\Debug
.\program1.exe
.\program2.exe
\end{lstlisting}
\begin{center}\small Листинг 10 - Команды сборки и запуска в Windows\end{center}

При использовании генератора с одной конфигурацией файлы могут находиться непосредственно в каталоге \texttt{build-win}. Запуск второй программы выполняется из каталога с DLL, поскольку в менеджере используются относительные пути.

Работоспособность обеих программ в Windows проверена. Все проведённые тесты пройдены успешно. Таким образом, функциональная проверка подтвердила работоспособность реализации на обеих операционных системах.

\section{Выводы}
В ходе лабораторной работы изучено создание динамических библиотек с общим интерфейсом и различными реализациями функций. Для контрактов \texttt{PrimeCount()} и \texttt{Square()} реализованы алгоритмы подсчёта простых чисел и вычисления площади двух фигур. Подготовлены две программы: одна использует библиотеку, подключённую при линковке, другая загружает библиотеки во время исполнения и переключает их по команде пользователя.

Сравнение двух подходов показало, что линковка с динамической библиотекой при сборке делает код программы проще: функции вызываются непосредственно, а загрузкой зависимости занимается система. Явная загрузка через \texttt{dlopen()} и \texttt{dlsym()} в Linux либо \texttt{LoadLibraryA()} и \texttt{GetProcAddress()} в Windows требует дополнительного кода для управления дескрипторами, указателями на функции и ошибками, но позволяет менять реализацию без перезапуска программы и её повторной линковки. В обоих случаях используются именно динамические библиотеки; различается способ получения доступа к их функциям.

Кроссплатформенность обеспечивается разделением общего кода и платформозависимых операций. Благодаря интерфейсу \texttt{os\_api.h} менеджер библиотек не зависит от конкретных функций ОС, а CMake выбирает нужную реализацию при сборке. Различия между платформами сводятся к механизму загрузки, типам дескрипторов, обработке результатов системных функций и формату библиотек.

Консольные тесты в Linux продемонстрировали работу обеих программ и переключение формулы площади в программе № 2. Тестирование в Windows также завершилось успешно. Результаты подтверждают работоспособность кроссплатформенной реализации на обеих операционных системах.
\end{document}
